\documentclass[conference]{IEEEtran}

\usepackage[utf8]{inputenc}
\usepackage[T1]{fontenc}
\usepackage[french,english]{babel}
\usepackage{cite}
\usepackage{amsmath,amssymb}
\usepackage{graphicx}
\usepackage{booktabs}
\usepackage{url}
\usepackage{xcolor}
\usepackage{microtype}
\usepackage{tikz}
\usetikzlibrary{shapes.geometric, arrows.meta, positioning, fit, backgrounds}

% ── TikZ styles ──────────────────────────────────────────────────────────────
\tikzset{
  block/.style={rectangle, rounded corners=3pt, draw=black!70, fill=blue!8,
                text width=2.6cm, align=center, font=\footnotesize, minimum height=0.7cm,
                inner sep=4pt},
  blockred/.style={block, fill=red!8},
  blockgreen/.style={block, fill=green!8},
  blockorange/.style={block, fill=orange!10},
  blockgray/.style={block, fill=gray!10},
  decision/.style={diamond, draw=black!70, fill=yellow!15, aspect=2.2,
                   font=\scriptsize, align=center, inner sep=2pt},
  arrow/.style={-{Stealth[length=5pt]}, thick, black!70},
  groupbox/.style={draw=black!30, dashed, rounded corners=5pt, fill=black!3,
                   inner sep=6pt},
}

\begin{document}

\title{Syst\`{e}me d'Aide au Diagnostic Radiologique par Segmentation S\'{e}mantique Multi-Organe et G\'{e}n\'{e}ration Automatique de Comptes Rendus en Fran\c{c}ais~: Application aux Donn\'{e}es Cliniques du Centre de Radiologie \'{E}milie, Gabon}

\author{
  \IEEEauthorblockN{Alaaeddine Bouchamla}
  \IEEEauthorblockA{ENISO, Universit\'{e} de Sousse\\Sousse, Tunisie\\
  alaaeddine.bouchamla@eniso.u-sousse.tn}
  \and
  \IEEEauthorblockN{Mehrez Abdellaoui}
  \IEEEauthorblockA{ENISO, Universit\'{e} de Sousse\\Sousse, Tunisie}
  \and
  \IEEEauthorblockN{Dahas Jalel \and Luis Felipe}
  \IEEEauthorblockA{Centre de Radiologie \'{E}milie (CRE)\\Libreville, Gabon}
}

\maketitle

% ── Abstract ─────────────────────────────────────────────────────────────────
\begin{abstract}
We present an end-to-end AI-assisted radiology reporting system combining
multi-organ semantic segmentation with automated French-language report
generation, deployed as a clinical web application for the Centre de Radiologie
\'{E}milie (CRE) in Libreville, Gabon. The system ingests DICOM CT and MRI
volumes, performs modality-aware segmentation using pretrained MONAI bundles
(104-organ CT segmentation via TotalSegmentator \cite{wasserthal2023}; 133-structure
brain MRI segmentation via UNesT \cite{yu2023}), applies a CRE-specific fine-tuning
layer trained on MedSAM \cite{ma2024} pseudo-labels from real patient data, and
generates draft \textit{comptes rendus} in French using a retrieval-augmented
generation (RAG) pipeline grounded in real institutional reports. A radiologist
reviews, edits, and validates each AI draft before PDF generation, ensuring
clinical safety. Validated reports are indexed back into the knowledge base,
implementing continual learning. To our knowledge, this is the first French-language
AI radiology reporting system built on real clinical data from sub-Saharan Africa.
\end{abstract}

\begin{IEEEkeywords}
semantic segmentation, MONAI, transfer learning, radiology report generation,
RAG, DICOM, French NLP, medical AI, sub-Saharan Africa
\end{IEEEkeywords}

% ── I. Introduction ───────────────────────────────────────────────────────────
\section{Introduction}

Radiology departments in sub-Saharan African countries face a dual challenge:
growing demand for imaging studies and a critical shortage of trained radiologists
\cite{osei2019}. The Centre de Radiologie \'{E}milie (CRE) in Libreville, Gabon
performs hundreds of CT and MRI examinations monthly, each requiring a structured
written \textit{compte rendu} produced manually --- a process subject to
inter-observer variability that creates reporting backlogs delaying patient care.

Deep learning has achieved near-expert performance in medical image segmentation,
with U-Net \cite{ronneberger2015} and its variants as the dominant paradigm.
Simultaneously, large language models enable automated clinical text generation
\cite{singhal2023}. Combining these into a single pipeline from raw DICOM input
to a validated radiology report has received limited attention for French-language
contexts and has not, to our knowledge, been applied to real African clinical data.

This work makes four contributions:
\begin{enumerate}
  \item A \textbf{modality-aware segmentation pipeline} using pretrained MONAI
        bundles at no additional compute cost, covering the full CRE case mix.
  \item A \textbf{CRE-specific fine-tuning layer} trained on pseudo-labels
        automatically generated by MedSAM --- requiring no manual annotation.
  \item An \textbf{automated French report generation system} using RAG from
        a growing corpus of real CRE \textit{comptes rendus}, grounded in
        segmentation findings.
  \item A \textbf{deployable human-in-the-loop clinical tool} with doctor
        review, patient portal, and continual learning loop.
\end{enumerate}

% ── II. Related Work ──────────────────────────────────────────────────────────
\section{Related Work}

\textbf{Medical image segmentation.}
The U-Net architecture \cite{ronneberger2015} introduced encoder-decoder design
with skip connections that remains dominant. TotalSegmentator \cite{wasserthal2023}
(SegResNet, 1{,}204 CT scans, 104 organs) and UNesT \cite{yu2023} (hierarchical
transformer, 133 brain structures) represent the current state of the art for
pretrained clinical segmentation. nnU-Net \cite{isensee2021} provides
self-configuring training. MedSAM \cite{ma2024}, a fine-tuned Segment Anything
Model, enables prompt-based zero-shot medical segmentation.

\textbf{Automated radiology report generation.}
Early systems relied on template filling \cite{shin2016}. Recent approaches ---
BioViL-T \cite{bannur2023} and LLaVA-Med --- use vision-language pretraining.
Most target English chest X-ray reports; French multi-organ reporting is
unexplored.

\textbf{RAG for medical applications.}
Retrieval-Augmented Generation \cite{lewis2020} improves factual grounding by
conditioning generation on retrieved evidence. Its application to radiology
report generation --- grounding text in real prior reports --- is novel.

\textbf{AI in sub-Saharan African healthcare.}
Published AI applications to African clinical data remain sparse \cite{osei2019}.
This work contributes a real-world dataset and deployed system from Gabon.

% ── III. System Architecture ──────────────────────────────────────────────────
\section{System Architecture}

\subsection{Pipeline Overview}

Fig.~\ref{fig:pipeline} shows the complete system pipeline. A FastAPI backend
exposes REST endpoints consumed by a Next.js web frontend. SQLite stores patient
records, study metadata, segmentation results, and reports. ChromaDB indexes all
\textit{comptes rendus} for RAG retrieval.

% ── Full-width pipeline figure ────────────────────────────────────────────────
\begin{figure*}[t]
\centering
\begin{tikzpicture}[node distance=0.45cm and 0.35cm]

  % ── Row 1: input ──
  \node[block] (dicom) {DICOM ZIP\\Upload};

  % ── Row 2: preprocess ──
  \node[block, below=of dicom] (pre) {Preprocessing\\{\scriptsize pydicom · pyjpegls}\\{\scriptsize HU window · 512\texttimes512}};

  % ── Row 3: router ──
  \node[decision, below=of pre, yshift=-0.1cm] (router) {Modality\\Router};

  % ── Row 4: three branches ──
  \node[block, below left=0.6cm and 1.8cm of router] (ct)
        {\textbf{Stage 1A}\\TotalSegmentator\\CT · 104 organs};
  \node[block, below=0.6cm of router] (brain)
        {\textbf{Stage 1B}\\UNesT\\Brain MRI\\133 structures};
  \node[blockred, below right=0.6cm and 1.8cm of router] (msk)
        {MedSAM\\Prompted\\MSK MRI\\{\scriptsize (limitation)}};

  % ── Row 5: merge → Stage 2 ──
  \node[blockorange, below=1.8cm of router] (stage2)
        {\textbf{Stage 2}\\CRE U-Net\\Fine-tuned on\\pseudo-labels};

  % ── Row 6: clinical layer ──
  \node[block, below=of stage2] (clinical)
        {\textbf{Clinical Layer}\\Volume (mL) · Mean HU\\Anomaly flags\\Overlay PNG};

  % ── Row 7: split into RAG + findings ──
  \node[block, below left=0.6cm and 1.5cm of clinical] (rag)
        {RAG Retrieval\\ChromaDB · k=5\\Multilingual\\embeddings};
  \node[block, below right=0.6cm and 1.5cm of clinical] (findings)
        {Findings JSON\\structure · size\\location · HU};

  % ── Row 8: LLaVA ──
  \node[block, below=1.7cm of clinical] (llava)
        {LLaVA-7B\\Ollama · local\\French medical\\R\'{e}sultat + Conclusion};

  % ── Row 9: doctor ──
  \node[blockgreen, below=of llava] (doctor)
        {Radiologist Review\\Web UI · Edit · Sign\\\textit{``BROUILLON IA''}};

  % ── Row 10: two outputs ──
  \node[block, below left=0.6cm and 1.5cm of doctor] (pdf)
        {PDF Generation\\ReportLab\\CRE branding\\Watermark};
  \node[blockgray, below right=0.6cm and 1.5cm of doctor] (cl)
        {Continual Learning\\Validated report\\indexed \textrightarrow{} ChromaDB\\RAG improves};

  % ── Arrows: main flow ──
  \draw[arrow] (dicom) -- (pre);
  \draw[arrow] (pre) -- (router);
  \draw[arrow] (router) -- node[left,font=\tiny]{CT} (ct);
  \draw[arrow] (router) -- node[right,font=\tiny]{Brain} (brain);
  \draw[arrow] (router) -- node[right,font=\tiny]{MSK} (msk);
  \draw[arrow] (ct) -- (stage2);
  \draw[arrow] (brain) -- (stage2);
  \draw[arrow] (msk) -- (stage2);
  \draw[arrow] (stage2) -- (clinical);
  \draw[arrow] (clinical) -- (rag);
  \draw[arrow] (clinical) -- (findings);
  \draw[arrow] (rag) -- (llava);
  \draw[arrow] (findings) -- (llava);
  \draw[arrow] (llava) -- (doctor);
  \draw[arrow] (doctor) -- (pdf);
  \draw[arrow] (doctor) -- (cl);

  % ── Group boxes ──
  \begin{scope}[on background layer]
    \node[groupbox, fit=(ct)(brain)(msk)(stage2),
          label={[font=\scriptsize,gray]above:Segmentation Pipeline}] {};
    \node[groupbox, fit=(rag)(findings)(llava),
          label={[font=\scriptsize,gray]above:Report Generation}] {};
  \end{scope}

\end{tikzpicture}
\caption{End-to-end pipeline of the AI-assisted radiology reporting system.
  The modality router dispatches each volume to the appropriate pretrained
  MONAI bundle (Stage~1), followed by CRE-specific fine-tuning (Stage~2),
  per-organ clinical metric extraction, RAG-grounded French report generation,
  and human-in-the-loop validation with continual learning.}
\label{fig:pipeline}
\end{figure*}

\subsection{Technology Stack}

\begin{table}[h]
\caption{Technology stack}
\label{tab:stack}
\centering
\renewcommand{\arraystretch}{1.1}
\begin{tabular}{ll}
\toprule
\textbf{Layer} & \textbf{Technology} \\
\midrule
Backend & FastAPI (Python 3.10+), async SQLAlchemy \\
Frontend & Next.js 14, TypeScript, Tailwind CSS \\
Segmentation Stage 1 & MONAI pretrained bundles (PyTorch) \\
Segmentation Stage 2 & Custom 2D U-Net, PyTorch \\
Pseudo-labeling & MedSAM ViT-B (HuggingFace) \\
RAG & ChromaDB, sentence-transformers \\
Report generation & Ollama + LLaVA-7B (local) \\
PDF & ReportLab \\
Training & Kaggle T4 GPU (free tier) \\
\bottomrule
\end{tabular}
\end{table}

% ── IV. Methodology ───────────────────────────────────────────────────────────
\section{Methodology}

\subsection{Stage 1 --- Pretrained Multi-Organ Segmentation}

Rather than training from scratch, we leverage two MONAI pretrained bundles
covering the full CRE clinical case mix at no additional compute cost.

\textbf{CT volumes} are routed to \texttt{wholeBody\_ct\_segmentation}
\cite{wasserthal2023} --- a SegResNet architecture trained on the TotalSegmentator
dataset (1{,}204 clinical CT scans, 104 anatomical structures expert-annotated).
We use the 3.0\,mm low-resolution variant, which fits a free T4 GPU. This single
model covers all CRE CT protocols: TDM Thoracique, TDM TAP, TDM
Abdomino-Pelvien, and generic CT.

\textbf{Brain MRI volumes} are routed to
\texttt{wholeBrainSeg\_Large\_UNEST\_segmentation} \cite{yu2023} --- a UNesT
architecture (hierarchical transformer + 3D U-Net decoder) trained on 133 brain
structures from T1-weighted MRI. This covers IRM C\'{e}r\'{e}bral cases.

\textbf{MSK MRI} (knee, shoulder, elbow, cervical spine) has no comparable
pretrained model publicly available. These cases are processed via MedSAM
prompted segmentation, documented as a system limitation.

Modality detection is performed automatically from DICOM metadata
(\texttt{Modality}, \texttt{BodyPartExamined} tags) with an HU
intensity-based heuristic fallback.

\subsection{Stage 2 --- CRE-Specific Fine-Tuning}

Stage~1 models are not trained on CRE data. To adapt segmentation to the
specific characteristics of CRE's scanner hardware and patient population, we
fine-tune a 2D U-Net on real CRE patient data using automatic pseudo-labels.

\textbf{Pseudo-labeling via MedSAM.}
For each patient, representative DICOM slices are extracted at uniform intervals.
Automatic bounding-box prompts are generated by Otsu intensity thresholding
followed by connected-component analysis (top-$N$ largest components).
MedSAM (ViT-B, \texttt{flaviagiammarino/medsam-vit-base}) produces binary
segmentation masks without any manual drawing. Slice/mask pairs are saved as
NumPy arrays with a manifest JSON.

\textbf{Fine-tuning setup.}
2D U-Net (DoubleConv blocks, skip connections),
$\text{in\_channels}=1$, $\text{num\_classes}=2$ (binary),
features $= [64, 128, 256, 512]$.
Loss: combined Dice + CrossEntropy ($\alpha=0.5$).
Optimizer: AdamW, $\text{lr}=5\!\times\!10^{-5}$, weight decay $10^{-5}$.
Scheduler: CosineAnnealingLR, $T_{\max}=25$.
Training: 25~epochs, checkpoint every~3~epochs, early stopping patience~6.
Encoder/decoder weights matching Stage~1 shape are transferred; only the
input convolution (4ch$\!\to\!$1ch) and output head ($N$-class$\!\to\!$2-class)
are freshly trained.

\subsection{Clinical Application Layer}

Segmentation masks are post-processed to extract per-organ metrics:
\textbf{volume} (mL = voxel count $\times$ voxel spacing$^3$),
\textbf{mean HU density},
\textbf{presence/absence flag},
and \textbf{anomaly flags} (organ volumes compared against published normal
ranges \cite{molina2021}).
A convolutional autoencoder provides an image-level anomaly score (MSE
reconstruction error; slices exceeding $\mu + 2\sigma$ are flagged).

\subsection{RAG-Based French Report Generation}

Structured findings are encoded as a query. ChromaDB retrieves the $k=5$ most
semantically similar \textit{comptes rendus} from the institutional corpus using
sentence-transformer embeddings (\textit{paraphrase-multilingual-MiniLM-L12-v2},
384-dim, multilingual, French-optimised). Findings and retrieved examples are
provided as context to Ollama (LLaVA-7B, local inference). The model generates
the \textit{R\'{e}sultat} and \textit{Conclusion} sections in French. A
rule-based template fallback activates when Ollama is unavailable.

\textbf{Key design principle.} Segmentation output directly feeds report text.
A detected liver volume of 2{,}100\,mL generates: \textit{``Foie augment\'{e} de
volume (2\,100\,mL), compatible avec une h\'{e}patomégalie --- \`{a} corr\'{e}ler
cliniquement.''} This coupling between visual analysis and language distinguishes
our approach from systems where segmentation and text generation are independent.

\subsection{Human-in-the-Loop Validation and Continual Learning}

The AI draft is watermarked \textit{``BROUILLON IA --- EN ATTENTE DE VALIDATION''}
throughout the UI and on every PDF page. The radiologist reviews the segmentation
overlay and editable draft side-by-side, makes corrections, and signs. The
validated \textit{compte rendu} is indexed back into ChromaDB via upsert ---
implementing continual learning without retraining.

% ── V. Experiments ────────────────────────────────────────────────────────────
\section{Experiments and Results}

\subsection{Stage 2 Segmentation Metrics}

\begin{table}[h]
\caption{Stage 2 segmentation metrics (CRE pseudo-label validation set)}
\label{tab:seg}
\centering
\renewcommand{\arraystretch}{1.1}
\begin{tabular}{lc}
\toprule
\textbf{Metric} & \textbf{Value} \\
\midrule
IoU (mean)  & {[FILL]} \\
Dice (mean) & {[FILL]} \\
Precision   & {[FILL]} \\
Recall      & {[FILL]} \\
\bottomrule
\multicolumn{2}{l}{\scriptsize Stage~1 MONAI models are evaluated by their}\\
\multicolumn{2}{l}{\scriptsize original authors \cite{wasserthal2023,yu2023}.}\\
\end{tabular}
\end{table}

\subsection{Report Generation Quality}

Dr.~Dahas Jalel evaluated [FILL] AI-generated \textit{comptes rendus} on five
criteria (scale 1--5):

\begin{table}[h]
\caption{Clinical evaluation by radiologist (Dr. Dahas Jalel)}
\label{tab:eval}
\centering
\renewcommand{\arraystretch}{1.1}
\begin{tabular}{lc}
\toprule
\textbf{Crit\`{e}re} & \textbf{Score moyen} \\
\midrule
Terminologie m\'{e}dicale fran\c{c}aise & {[FILL]} \\
Pertinence clinique des r\'{e}sultats   & {[FILL]} \\
Compl\'{e}tude du compte rendu          & {[FILL]} \\
Respect du format CRE                   & {[FILL]} \\
Utilit\'{e} pour le radiologue          & {[FILL]} \\
\midrule
\textbf{Score global}                   & \textbf{[FILL] / 5} \\
\bottomrule
\end{tabular}
\end{table}

BLEU-4: [FILL] \quad ROUGE-L: [FILL] \quad ($n=$[FILL] held-out CT reports)

\subsection{System Performance}

\begin{table}[h]
\caption{End-to-end inference time}
\label{tab:perf}
\centering
\renewcommand{\arraystretch}{1.1}
\begin{tabular}{lcc}
\toprule
\textbf{Component} & \textbf{CPU} & \textbf{GPU (GTX 1650)} \\
\midrule
DICOM decoding (610 slices) & $\sim$3\,s  & $\sim$1\,s \\
Stage 1 segmentation (5 sl.) & $\sim$12\,s & $\sim$2\,s \\
Per-organ metrics            & $\sim$0.5\,s & $\sim$0.5\,s \\
RAG retrieval ($k=5$)        & $\sim$1\,s  & $\sim$1\,s \\
LLaVA report generation      & $\sim$45\,s & $\sim$45\,s \\
PDF generation               & $\sim$0.5\,s & $\sim$0.5\,s \\
\midrule
\textbf{Total pipeline}      & $\sim$62\,s & $\sim$50\,s \\
\bottomrule
\end{tabular}
\end{table}

% ── VI. Discussion ────────────────────────────────────────────────────────────
\section{Discussion}

\textbf{Clinical value.}
The system reduces \textit{compte rendu} production time from 10--20\,min of
manual writing to a 2--5\,min review-and-edit task. For normal examinations ---
a large proportion of CRE studies --- the AI draft is often immediately usable
with minor corrections. The RAG component ensures institutional style
consistency.

\textbf{Novelty.}
Three aspects are, to our knowledge, novel: (1)~use of MONAI pretrained bundles
in a deployed clinical system with modality-aware routing; (2)~MedSAM
pseudo-labeling as a zero-annotation strategy for CRE-specific adaptation;
(3)~the combination of segmentation $\to$ volumetric metrics $\to$ RAG $\to$
French report $\to$ continual learning in a single deployed application on real
sub-Saharan African clinical data.

\textbf{Limitations.}
MSK MRI (knee, shoulder, elbow, spine) lacks a comparable pretrained
segmentation model. Stage~2 training data is currently limited to [FILL]
patients; performance will improve as the CRE dataset grows. LLaVA may
occasionally hallucinate --- mandatory radiologist validation mitigates this risk.

\textbf{Future work.}
Indexing the full CRE corpus to maximise RAG quality; LoRA fine-tuning of a
French medical LLM (CamemBERT-Bio) on the GTX~1650; formal clinical trial
comparing AI-assisted vs.\ unassisted reporting; production deployment at CRE.

% ── VII. Conclusion ───────────────────────────────────────────────────────────
\section{Conclusion}

We presented an end-to-end AI-assisted radiology reporting system for the Centre
de Radiologie \'{E}milie, Gabon --- combining modality-aware multi-organ
segmentation (MONAI pretrained bundles + CRE fine-tuning), volumetric clinical
metrics, RAG-based French report generation, and human-in-the-loop validation.
The system is clinically safe by design and implements continual learning through
doctor-validated report indexing. It represents, to our knowledge, the first such
system built on real clinical data from sub-Saharan Africa.

% ── References ────────────────────────────────────────────────────────────────
\bibliographystyle{IEEEtran}
\begin{thebibliography}{11}

\bibitem{wasserthal2023}
J.~Wasserthal \textit{et al.}, ``TotalSegmentator: Robust Segmentation of 104
Anatomical Structures in CT Images,'' \textit{Radiology: Artificial Intelligence},
vol.~5, no.~5, 2023.

\bibitem{yu2023}
T.~Yu \textit{et al.}, ``UNesT: Local Spatial Representation Learning with
Hierarchical Transformer for Efficient Medical Segmentation,''
\textit{Medical Image Analysis}, vol.~90, 2023.

\bibitem{ma2024}
J.~Ma \textit{et al.}, ``Segment Anything in Medical Images,''
\textit{Nature Communications}, vol.~15, no.~1, 2024.

\bibitem{osei2019}
A.~Osei-Bonsu \textit{et al.}, ``Radiology workforce in Africa: a situational
analysis,'' \textit{Pan African Medical Journal}, vol.~34, 2019.

\bibitem{ronneberger2015}
O.~Ronneberger, P.~Fischer, and T.~Brox, ``U-Net: Convolutional Networks for
Biomedical Image Segmentation,'' in \textit{MICCAI}, 2015, pp.~234--241.

\bibitem{singhal2023}
K.~Singhal \textit{et al.}, ``Large Language Models Encode Clinical Knowledge,''
\textit{Nature}, vol.~620, pp.~172--180, 2023.

\bibitem{isensee2021}
F.~Isensee \textit{et al.}, ``nnU-Net: a self-configuring method for deep
learning-based biomedical image segmentation,'' \textit{Nature Methods},
vol.~18, pp.~203--211, 2021.

\bibitem{shin2016}
H.~Shin \textit{et al.}, ``Learning to Read Chest X-Rays: Recurrent Neural
Cascade Model for Automated Image Annotation,'' in \textit{CVPR}, 2016.

\bibitem{bannur2023}
S.~Bannur \textit{et al.}, ``Learning to Exploit Temporal Structure for
Biomedical Vision--Language Processing,'' in \textit{CVPR}, 2023.

\bibitem{lewis2020}
P.~Lewis \textit{et al.}, ``Retrieval-Augmented Generation for
Knowledge-Intensive NLP Tasks,'' in \textit{NeurIPS}, 2020.

\bibitem{molina2021}
D.~Molina \textit{et al.}, ``Normal organ volumes and reference doses in
radiotherapy planning,'' \textit{Radiotherapy and Oncology}, vol.~155,
pp.~225--233, 2021.

\end{thebibliography}

\end{document}
